\chapterimage{chapter_head_2.pdf}
\chapter{Tries and Hashing}

\section{Theoretical Core}

\subsection{Tries}
\begin{definition}[Trie]
A multiway tree (often called a prefix tree) used to store a list of strings. Nodes represent characters (often implied by the index in a child pointer array), and the path from the root to a node defines a prefix or word.
\end{definition}

\begin{itemize}
    \item \textbf{Node Structure:} Contains an array of 26 pointers (for 'a'-'z') and a boolean flag \texttt{is\_end\_}\allowbreak\texttt{of\_word}.
    \item \textbf{Complexity:} $O(k)$ for Insert, Search, and Remove, where $k$ is the length of the string.
    \item \textbf{Space Trade-off:} Fast prefix searches at the cost of high memory usage due to many NULL pointers.
\end{itemize}

\subsection{Hash Tables}
\begin{definition}[Hash Table]
An implementation of the Table ADT that stores (key, value) pairs and allows for $O(1)$ expected time retrieval using a hash function.
\end{definition}

\begin{itemize}
    \item \textbf{Hash Function ($h(k)$):} Maps a key $k$ to a slot index. Common method: $h(k) = k \pmod m$.
    \item \textbf{Load Factor ($\alpha$):} Ratio of elements $n$ to slots $m$ ($\alpha = n/m$).
    \item \textbf{Collision Resolution:}
    \begin{itemize}
        \item \textbf{Separate Chaining:} Each slot points to a linked list.
        \item \textbf{Open Addressing:} Probe for the next available slot.
    \end{itemize}
\end{itemize}

\section{Algorithmic Tracing}

\begin{example}[Hashing with Probing]
Table size $m = 7$, $h(k) = k \pmod 7$. Insert keys: $A(6), B(6), C(5), D(5)$.

\textbf{Linear Probing ($h(k, i) = (h'(k) + i) \pmod 7$):}
\begin{enumerate}
    \item $A$: $6 \pmod 7 = 6$. Insert at index 6.
    \item $B$: $6 \pmod 7 = 6$ (Coll). Probe $i=1$: $(6+1) \pmod 7 = 0$. Insert at 0.
    \item $C$: $5 \pmod 7 = 5$. Insert at index 5.
    \item $D$: $5 \pmod 7 = 5$ (Coll). Probe $i=1$: $(5+1) \pmod 7 = 6$ (Coll). Probe $i=2, 3$: Insert at 1.
\end{enumerate}

\textbf{Quadratic Probing ($h(k, i) = (h'(k) + i^2) \pmod 7$):}
\begin{enumerate}
    \item $A, B, C$: Same as linear (at indices 6, 0, 5).
    \item $D$: $5 \pmod 7 = 5$ (Coll). Probe $i=1$: $(5+1^2) \pmod 7 = 6$ (Coll). Probe $i=2$: $(5+2^2) \pmod 7 = 2$. Insert at 2.
\end{enumerate}
\end{example}

\section{Code Implementation}

\begin{lstlisting}[language=C]
#define MAX_CHILD_COUNT 26

typedef struct TrieNode_s {
    struct TrieNode_s *children[MAX_CHILD_COUNT];
    int is_end_of_word;
} TrieNode;

// Simple String Hash Function
int string_to_int(char *str) {
    int sum = 0;
    for (int i = 0; str[i] != '\0'; i++) {
        sum += str[i];
    }
    return sum;
}

int hash_function(int key, int m) {
    return key % m;
}
\end{lstlisting}

\section{Skills \& Pitfalls}

\subsection{Must-Know Skills}
\begin{itemize}
    \item \textbf{Prefix Matching:} Tries are superior for finding all words starting with "pre" by traversing to the "e" node and exploring the subtree.
    \item \textbf{Table Sizing:} For Quadratic Probing, ensure the table size $m$ is a prime number and the load factor $\alpha < 0.5$ to guarantee finding an empty slot.
\end{itemize}

\subsection{Common Mistakes (Pitfalls)}
\begin{remark}
\begin{itemize}
    \item \textbf{Clustering:} Linear probing suffers from \textbf{primary clustering}, where long runs of occupied slots increase search time. Quadratic probing mitigates this but can cause \textbf{secondary clustering}.
    \item \textbf{Tombstones:} When deleting from an open-addressing table, use a \texttt{DELETED} marker. Clearing the slot to \texttt{NULL} breaks the probe sequence for other keys.
    \item \textbf{Trie Memory Management:} Always use a post-order traversal to \texttt{free()} trie nodes to avoid massive memory leaks.
\end{itemize}
\end{remark}
